% config.tex
\ProvidesFile{config.tex}[2026 NSFC Common Config]

% =========================================================
% 模块 1：基础宏包引入 
% =========================================================
\usepackage[english]{babel}
\usepackage{xeCJK}
\usepackage{xeCJKfntef}
\usepackage{fontspec} 
\usepackage{pifont}
\usepackage{graphicx}
\usepackage{ulem}
\usepackage{setspace}
\usepackage{enumitem} % 列表控制
\usepackage{amsmath}
\usepackage{mdframed}
\usepackage{xcolor}
\usepackage{amsfonts}
\usepackage{booktabs}
\usepackage[subrefformat=parens,labelformat=parens]{subfig}
\usepackage{hyperref}
\usepackage{titlesec} 
\usepackage{tcolorbox}
\usepackage{amssymb}  % 提供 \blacksquare 等数学符号
\usepackage{bm}       % 提供数学粗体

% =========================================================
% 模块 2：读取外部用户配置 (置于前方，确保全局变量生效)
% =========================================================
% =========================================================
%  NSFC Style Configuration (风格配置文件)
%  说明：此文件仅用于控制排版风格开关，核心逻辑在 config.tex
% =========================================================

% ---------------------------------------------------------
%  开关 1：正文字体设置
% ---------------------------------------------------------
% true  = 正文强制使用【楷体】
% false = 正文使用标准【宋体】
\newif\ifUseKaitiBody
\UseKaitiBodyfalse  % <--- 修改这里：true 或 false

% ---------------------------------------------------------
%  开关 2：标题字体选择
% ---------------------------------------------------------
% 定义各级标题（一级、二级）使用的字体
% 常用选项：\kaishu (楷体), \heiti (黑体), \songti (宋体)
\newcommand{\HeadingFont}{\mysong} 

% ---------------------------------------------------------
%  开关 3：标题编号风格
% ---------------------------------------------------------
% 请输入数字选择一种风格：
% 1 = 标准简洁: 1  / 1.1 / 1.1.1 
% 2 = Word风格: 1. / 1.1 / 1.1.1 (一级标题带个点)
% 3 = 中文风格: (一) / 1. / (1)  (公文风)
\newcommand{\NumberingStyle}{2} % <--- 修改这里：1, 2, 或 3

% ---------------------------------------------------------
%  开关 4：交叉引用与链接颜色
% ---------------------------------------------------------
% 定义参考文献标号、目录、公式引用的颜色。
% 打印版（黑白纸质）强烈建议用 black。电子版可用 blue, red, green, darkgray 等
\newcommand{\ThemeColor}{blue} % <--- 修改这里：输入你想要的颜色单词


% ---------------------------------------------------------
%  开关 5：全局布局与间距调控
% ---------------------------------------------------------
% 设置说明：
% 1. 段间距 (UserParskip): 建议值 0.3ex 到 0.8ex。若篇幅吃紧可设为 0pt。
\newcommand{\UserParskip}{0.5ex} % <--- 修改这里：段间距数值 (如 0.5ex, 3pt)

% 2. 列表模式 (UserListMode): 
%    0 = 极致紧凑: 列表项内外完全不留白 (nosep)，最省空间。
%    1 = 宽松美观: 列表间距随段间距自动按比例缩放，呼吸感强。
\newcommand{\UserListMode}{0}    % <--- 修改这里：0 或 1

%-----------------------------------------------
% 自定义命令（自定义颜色加粗）
%-----------------------------------------------
\definecolor{myemphasize}{RGB}{0, 0, 64} % 经典学术蓝，比纯黑更有质感#000640
% 定义自定义加粗命令 \mybf
\newcommand{\mybf}[1]{{\color{myemphasize}\textbf{#1}}}


% =========================================================
% 模块 3：基础格式与环境配置
% =========================================================
% --- 颜色与路径 ---
\definecolor{NsfcBlue}{RGB}{0,112,192}
\graphicspath{{figures/}}

% --- 参考文献设置 ---
\usepackage[square,numbers,sort&compress,super]{natbib}
\setlength{\bibsep}{1pt plus 0.3ex}
\pagestyle{empty}

% --- 中英文字体与间距设置 ---
\CJKsetecglue{\hskip 0.2em} % 中英文/数字间距
\renewcommand{\baselinestretch}{1.3} 

\setmainfont{Times New Roman}
\setCJKfamilyfont{mysong}[
    UprightFont = {fonts/simsun.ttc}, 
    % 针对粗体单独设置 Scale
    BoldFont    = {fonts/siyuansongti-Bold.otf},
    BoldFeatures = {Scale=0.96} 
]{}

% 创建快捷命令
\newcommand{\mysong}{\CJKfamily{mysong}}

\newfontfamily\KaiEN[Path=fonts/, AutoFakeBold=1.8]{simkai.ttf}
\setCJKfamilyfont{KaiCJK}[Path=fonts/, AutoFakeBold=1.8]{simkai.ttf}
\newcommand{\mykai}{\CJKfamily{KaiCJK}\KaiEN}

% --- 字号快捷命令 ---
\newcommand{\sanhao}{\zihao{3}}
\newcommand{\sihao}{\zihao{4}}
\newcommand{\xiaosihao}{\zihao{-4}}
\newcommand{\wuhao}{\zihao{5}}

% --- 图表标题中文化 ---
\addto\captionsenglish{
    \renewcommand{\refname}{\sihao \mykai \leftline{参考文献}}
    \renewcommand{\tablename}{表}
    \renewcommand{\figurename}{\mykai 图}
}

% --- 列表与代码块环境 ---
\setlist{nosep} % 清除列表多余间距

\newtcolorbox{codebox}{
    colback=gray!5,       
    colframe=gray!30,     
    boxrule=0.5pt,        
    arc=2pt,              
    left=4pt, top=2pt, bottom=2pt, right=4pt, 
    before skip=1ex, after skip=1ex 
}

% --- 超链接颜色设置 ---
\hypersetup{
    colorlinks = true,              
    citecolor  = \ThemeColor,       
    linkcolor  = \ThemeColor,       
    urlcolor   = \ThemeColor        
}

% =========================================================
% 模块 4：标题排版与自定义宏 (NSFC 专用)
% =========================================================

% --- 1. 动态间距长度定义 (底层 Hack 核心) ---
% 预设基础间距
\newlength{\SecTopBase}
\setlength{\SecTopBase}{2.0ex plus .2ex minus .2ex}
\newlength{\SubSecTopBase}
\setlength{\SubSecTopBase}{1.2ex plus .1ex minus .2ex}
\newlength{\SubSubSecTopBase}
\setlength{\SubSubSecTopBase}{0.8ex plus .1ex minus .1ex}

% 定义当前生效的动态间距 (初始等于基础间距)
\newlength{\SecTopCurrent}
\setlength{\SecTopCurrent}{\SecTopBase}
\newlength{\SubSecTopCurrent}
\setlength{\SubSecTopCurrent}{\SubSecTopBase}
\newlength{\SubSubSecTopCurrent}
\setlength{\SubSubSecTopCurrent}{\SubSubSecTopBase}

% --- NSFC 智能标题配置 ---
\newlength{\NsfcHeaderIndent}
\setlength{\NsfcHeaderIndent}{2em}

\newlength{\NsfcSubHeaderIndent}
\setlength{\NsfcSubHeaderIndent}{2em}

\newlength{\NsfcHeaderTopSkip}
\setlength{\NsfcHeaderTopSkip}{0.5em}

\newif\ifIsFirstSubHeader
\IsFirstSubHeadertrue

\newcounter{NsfcHeaderCounter}
\renewcommand{\theHsection}{\arabic{NsfcHeaderCounter}.\arabic{section}}

\newcommand{\NsfcHeader}[1]{%
    \stepcounter{NsfcHeaderCounter}%
    \setcounter{section}{0}%
    % 【关键】遇到全局大标题，把紧接在后面的所有子级标题顶部间距全部强行归零
    \global\setlength{\SecTopCurrent}{0pt}%
    \global\setlength{\SubSecTopCurrent}{0pt}%
    \par
    \vspace{\NsfcHeaderTopSkip}
    \noindent\hspace{\NsfcHeaderIndent}
    {\color{NsfcBlue}\sihao\mykai #1}
    \par
    \vspace{0.0em}
    \global\IsFirstSubHeadertrue
}

\newcommand{\NsfcSubHeader}[1]{%
    \par
    \ifIsFirstSubHeader
        \vspace{0.0em}
        \global\IsFirstSubHeaderfalse
    \else
        \vspace{0.2em}
    \fi
    \noindent\hspace{\NsfcSubHeaderIndent}
    {\sihao\color{NsfcBlue}\mykai\ziju{-0.00} #1}
    \vspace{0mm}
    % 【关键】遇到 NSFC 二级大标题，同样杀掉紧跟的正文标题间距
    \global\setlength{\SecTopCurrent}{0pt}%
    \global\setlength{\SubSecTopCurrent}{0pt}%
}

% --- 系统各级标题排版设置 (\section 等) ---
\ifcase\NumberingStyle
    % Case 0: (空)
\or % Case 1: 标准
    \renewcommand{\thesection}{\arabic{section}}
    \renewcommand{\thesubsection}{\arabic{section}.\arabic{subsection}}
    \renewcommand{\thesubsubsection}{\arabic{section}.\arabic{subsection}.\arabic{subsubsection}}
\or % Case 2: Word点号
    \renewcommand{\thesection}{\arabic{section}.}
    \renewcommand{\thesubsection}{\arabic{section}.\arabic{subsection}}
    \renewcommand{\thesubsubsection}{\arabic{section}.\arabic{subsection}.\arabic{subsubsection}}
\or % Case 3: 中文
    \renewcommand{\thesection}{（\chinese{section}）} 
    \renewcommand{\thesubsection}{\arabic{subsection}.}
    \renewcommand{\thesubsubsection}{（\arabic{subsubsection}）}
\fi

% --- 动态间距恢复机制 ---
% 逻辑：在排版标题主体之前（此时上方间距已经按照 0pt 结算完毕），将下一个同级标题的间距恢复为常态，并将子级标题的间距归零。
\titleformat{\section}{\zihao{4}\HeadingFont\bfseries}{\thesection}{0.5em}{%
    \global\setlength{\SecTopCurrent}{\SecTopBase}% 恢复下一个 \section 的间距
    \global\setlength{\SubSecTopCurrent}{0em}% 杀掉当前 \section 下第一个 \subsection 的间距
}
\titlespacing{\section}{0pt}{\SecTopCurrent}{1.0ex plus .2ex minus .2ex}

\titleformat{\subsection}{\zihao{4}\HeadingFont\bfseries}{\thesubsection}{0.5em}{%
    \global\setlength{\SubSecTopCurrent}{\SubSecTopBase}% 恢复下一个 \subsection 的间距
    \global\setlength{\SubSubSecTopCurrent}{0em}% 杀掉当前 \subsection 下第一个 \subsubsection 的间距
}
\titlespacing{\subsection}{0pt}{\SubSecTopCurrent}{0.5ex plus .1ex minus .2ex}

\titleformat{\subsubsection}{\zihao{-4}\HeadingFont\bfseries}{\thesubsubsection}{0.5em}{%
    \global\setlength{\SubSubSecTopCurrent}{\SubSubSecTopBase}% 恢复下一个 \subsubsection 的间距
}
\titlespacing{\subsubsection}{0pt}{\SubSubSecTopCurrent}{0.3ex plus .1ex minus .1ex}

% =========================================================
% 模块 5：核心文档状态注入 (合并全部 \AtBeginDocument)
% =========================================================

% 保留旧版废弃开关，防止外部主文档误调用导致直接报错
\newif\ifUseKaiBody
\UseKaiBodyfalse 

% 定义隐身模式开关
\newif\ifHideBody
\HideBodyfalse

% 统一管控正文字体、字号与隐身状态
\AtBeginDocument{
    \ifHideBody
        % 【隐身模式】: 字体变白，字号极小，行距压扁
        \color{white}
        \fontsize{1pt}{1pt}\selectfont
        \linespread{0.1}\selectfont
    \else
        % 【正常模式】: 强制正文小四号，并根据 style.tex 决定是否启用楷体
        \xiaosihao 
        \ifUseKaitiBody
            \ifdefined\mykai \mykai \else \kaishu \fi
        \else
            \ifdefined\mysong \mysong \else \songti \fi
        \fi
    \fi
}

% =========================================================
% 局部修复公式环境上下间距 (仅针对 equation 和 align 生效)
% =========================================================
\newcommand{\TightEquationSpacing}{%
    % 基础间距 0.5em，允许上下浮动 0.1em 到 0.2em 的弹性空间
    \setlength{\abovedisplayskip}{0.6em plus 0.2em minus 0.1em}%
    \setlength{\belowdisplayskip}{0.6em plus 0.2em minus 0.1em}%
    % 短行公式上方间距需进一步压缩，保持视觉紧凑
    \setlength{\abovedisplayshortskip}{0.2em plus 0.1em}%
    \setlength{\belowdisplayshortskip}{0.5em plus 0.2em minus 0.1em}%
}

\AtBeginEnvironment{equation}{\TightEquationSpacing}
\AtBeginEnvironment{align}{\TightEquationSpacing}
% 如果你还用了其他的 amsmath 环境，比如 gather, flalign 等，按同理追加即可

%--------------------------------------------------------
% 段间距
%--------------------------------------------------------
% 1. 声明长度变量
\newlength{\NsfcParskip}

% 2. 执行段间距设置 (利用 plus/minus 增加弹性)
\setlength{\NsfcParskip}{\UserParskip plus 0.2ex minus 0.1ex}
\setlength{\parskip}{\NsfcParskip}

% 3. 执行列表设置 (由于包含 if 逻辑，建议放在 AtBeginDocument 避错)
\AtBeginDocument{
    \ifnum\UserListMode=0
        \setlist{nosep} % 彻底无间距
    \else
        \setlist{
            partopsep=0pt,
            topsep=0.5\NsfcParskip,
            parsep=\NsfcParskip,
            itemsep=0.3\NsfcParskip
        }
    \fi
}